\documentclass[11pt]{article}
\usepackage[margin=1in]{geometry}
\usepackage{amsmath}
\usepackage{booktabs}
\usepackage{array}
\usepackage[hidelinks]{hyperref}
\usepackage[T1]{fontenc}
\usepackage{lmodern}

\title{Beyond Film Thickness: A Design-Only Research Plan for Correlative Geometric--Molecular--Force Metrology of Nanoscale Film Stability}
\author{Research Plan Author Artifact}
\date{30 August 2026}

\begin{document}
\maketitle

\noindent\textbf{Status:} PROPOSAL\_NO\_OBSERVED\_RESULTS\\
\textbf{Evidence status:} DESIGNED\_NOT\_EXECUTED\\
\textbf{Execution policy:} DESIGN\_ONLY\\
\textbf{Risk-review status:} BLOCKED\_BY\_RISK\_REVIEW

\begin{abstract}
Microscopic interfaces govern drainage, wetting, reaction, and transport, yet a film-thickness trace does not identify the molecular or mechanical causes of an interfacial transition. This research plan integrates a manual literature Survey, an Idea-stage causal hypothesis, and a design-only ExperimentDesign artifact into a single proposal for non-biological liquid--solid nanoscale thin-film interfaces. Its central proposition is conditional: within an expert-approved, geometry-matched comparison boundary, a temporally coherent relation among film geometry, an interface-specific molecular-state representation, and a local confined-interaction representation may distinguish chemistry-mediated disjoining-pressure control of film drainage or stability from a thickness-only explanation.

The contribution is an evidence architecture rather than an experimental result. It requires a calibrated geometric baseline, a molecular-state channel with explicit interpretation limits, a local interaction proxy with a perturbation audit, a declared transition outcome, and a common time/uncertainty representation. It compares a geometry-only account with single-added-channel and joint-channel accounts while retaining capillarity, optical-model assumptions, surface history, roughness, probe effects, observation-induced change, and temporal mismatch as live alternatives. No material system, chemical identity, instrument configuration, operating parameter, physical procedure, or observed result is supplied.
\end{abstract}

\noindent\textbf{Keywords:} interfacial metrology; nanoscale liquid films; sum-frequency generation; confined forces; disjoining pressure; causal discrimination; research design.

\section{Introduction}
Gas--liquid, liquid--solid, and solid--solid interfaces couple geometry, molecular organization, and transport across different spatial and temporal scales. Optical interference and ellipsometric methods provide nanometre-scale film geometry, but thickness is only one state variable. It does not uniquely identify molecular orientation, adsorption, local confined interactions, or the route by which a film drains, arrests, persists, or ruptures.

The literature supplies complementary component capabilities. Nonlinear optical spectroscopy can constrain interfacial molecular arrangement, while quantitative interpretation of a vibrational representation retains orientational and spectral assumptions~\cite{backus2020,wang2005}. Surface-force methods provide local confined-force and reactivity evidence but retain probe, roughness, and interaction-model limitations~\cite{dziadkowiec2019}. Imaging ellipsometry provides a non-contact geometric anchor for nanometre-scale liquid films~\cite{shoji2024}. These sources do not establish the combined causal relation proposed here.

This is a research plan, not an empirical report. It selects neither a material system nor an experimental procedure. The design is restricted to future expert-approved non-biological liquid--solid model interfaces.

\section{Research Objective and Causal Proposition}
The objective is to formulate a falsifiable, auditable plan for deciding whether a correlative geometric--molecular--force record adds explanatory information beyond a thickness-only account. The proposition is conditional: within a bounded non-biological liquid--solid thin-film system, a temporally coherent relation among geometry, an interface-specific vibrational state, and a confined-force proxy can distinguish chemistry-mediated disjoining-pressure control of drainage or stability from a thickness-only explanation.

For planning purposes, the candidate interfacial state at relative event time $t$ is represented by
\begin{equation}
\mathbf{X}(t)=\bigl[G(t),M(t),F(t),T(t),A(t)\bigr],
\label{eq:state}
\end{equation}
where $G$ is calibrated geometry, $M$ an interface-specific molecular-state representation, $F$ a local confined-interaction representation, $T$ a declared transition outcome, and $A$ the audit record for calibration, surface history, perturbation, and temporal registration. Equation~\eqref{eq:state} organizes future evidence; it does not claim that all variables are already measurable or causally complete.

The conceptual comparison is
\begin{equation}
\mathcal{M}_0:T\leftarrow G,\qquad \mathcal{M}_1:T\leftarrow(G,M,F,A),
\label{eq:comparison}
\end{equation}
where $\mathcal{M}_0$ is the geometry-only account and $\mathcal{M}_1$ is the joint account. No numerical threshold for preferring $\mathcal{M}_1$ is fixed. Any future criterion requires qualified approval of the model system, calibration traceability, uncertainty model, and estimand.

\section{Evidence Architecture}
\begin{table}[htbp]
\centering
\caption{Complementary evidence roles and retained limitations.}
\label{tab:evidence}
\begin{tabular}{p{0.18\linewidth}p{0.34\linewidth}p{0.36\linewidth}}
\toprule
Role & Design requirement & Limitation retained \\
\midrule
Geometry & Calibrated, time-indexed film thickness and/or morphology & Optical model and refractive-index assumptions \\
Molecular state & Interface-specific vibrational or chemical-state representation & Spectral, phase, and orientational interpretation limits \\
Local interaction & Confined-force or disjoining-pressure representation & Probe perturbation, roughness, and interaction-model error \\
Transition outcome & Approved definition of drainage, arrest, persistence, or rupture & Endpoint is system-dependent \\
Audit & Calibration, surface history, perturbation, and time-registration record & Completeness must be assessed before interpretation \\
\bottomrule
\end{tabular}
\end{table}

The proposed accounts are geometry-only ($G,T,A$), geometry plus molecular state ($G,M,T,A$), geometry plus local interaction ($G,F,T,A$), and the joint account ($G,M,F,T,A$). These are explanatory accounts, not physical sample recipes. A valid future comparison must use one declared comparability boundary rather than presume that different observation channels are automatically equivalent.

An audit-first rule applies. A causal interpretation becomes eligible only after the geometry-only baseline, cross-channel uncertainty, time basis, surface-history record, and perturbation assessment are documented. If a critical audit is absent or non-discriminating, the evidence is uninformative or invalid rather than provisionally mechanistic.

\section{Analysis, Robustness, and Reproducibility}
After expert approval of an estimand, a future analysis may compare the explanatory adequacy of the models in Equation~\eqref{eq:comparison} under the same registered comparability boundary. This manuscript specifies no statistic, effect magnitude, acceptance threshold, allocation sequence, repeat count, or operating condition.

Conceptual ablations remove the molecular-state representation or the local-interaction representation to test whether the interpretation collapses to geometry or remains underdetermined. Calibration, registration, surface-history, roughness, probe, and observation-induced alternatives remain explicit sensitivity cases. Reproducibility requires preservation of raw and derived channel records, calibration metadata, time-registration decisions, surface-history metadata, uncertainty assumptions, comparison-model versions, exclusions, and outcome-branch selection.

\section{Conditional Outcome Branches}
\begin{table}[htbp]
\centering
\caption{Permitted interpretations of future outcomes.}
\label{tab:outcomes}
\begin{tabular}{p{0.21\linewidth}p{0.33\linewidth}p{0.34\linewidth}}
\toprule
Branch & Future trigger & Permitted interpretation \\
\midrule
Supports mechanism & Joint explanatory information exceeds geometry while artifact alternatives are insufficient & Bounded evidence is consistent with the mechanism only in the documented model boundary \\
Partial or heterogeneous & The relation is informative only for a subset of states or audits & The relation may be conditional or incompletely identifiable \\
Null or contradictory & Geometry explains equally well, or added channels conflict & The proposed mechanism is not supported in the documented boundary \\
Uninformative or invalid & A critical audit is absent or non-discriminating & No causal conclusion is permitted \\
\bottomrule
\end{tabular}
\end{table}

\section{Responsible-Research Boundary}
The ExperimentDesign artifact is DESIGN\_ONLY and its risk gate is BLOCKED\_BY\_RISK\_REVIEW. No people, animals, plants, clinical populations, environmental samples, or restricted biological systems are included. No material identity, chemical formulation, apparatus configuration, quantity, operating parameter, or physical procedure is provided.

Any future physical route requires qualified confirmation of chemical safety, facility compatibility, model-system appropriateness, calibration traceability, and governance requirements. These decisions remain needs\_human\_input. This manuscript cannot authorize physical experimentation.

\section{Conclusion}
This plan shifts the target from thickness measurement to evidence-architecture design. It defines a finite scientific object, a geometry-only baseline, complementary molecular and local-interaction representations, a causal alternative set, and explicit no-information conditions. It does not assert that chemistry-mediated disjoining pressure has been established. Instead, it specifies the evidence conditions under which a future non-biological study could discriminate that account from a thickness-only explanation, or correctly conclude that the evidence is null, contradictory, or non-interpretable.

\begin{thebibliography}{9}
\bibitem{backus2020} E. H. G. Backus, J. Schaefer, and M. Bonn, ``Probing the Mineral-Water Interface with Nonlinear Optical Spectroscopy,'' \textit{Angewandte Chemie International Edition}, 2020. DOI: 10.1002/anie.202003085.
\bibitem{wang2005} H. Wang, W. Gan, R. Lu, Y. Rao, and B. Wu, ``Quantitative spectral and orientational analysis in surface sum frequency generation vibrational spectroscopy (SFG-VS),'' \textit{International Reviews in Physical Chemistry}, 2005. DOI: 10.1080/01442350500225894.
\bibitem{dziadkowiec2019} J. Dziadkowiec, B. Zareeipolgardani, D. K. Dysthe, and A. Røyne, ``Nucleation in confinement generates long-range repulsion between rough calcite surfaces,'' \textit{Scientific Reports}, 2019. DOI: 10.1038/s41598-019-45163-6.
\bibitem{shoji2024} E. Shoji, A. Hoshino, T. Biwa, et al., ``Superspreading Wetting of Nanofluid Droplet Laden with Highly Dispersed Nanoparticles,'' \textit{Langmuir}, 2024. DOI: 10.1021/acs.langmuir.4c03347.
\end{thebibliography}

\end{document}
